This thesis investigated data-parallel scanning strategies and
microarchitectural optimizations to accelerate the Aho-Corasick multi-pattern
matching algorithm on shared-memory multi-core CPUs, using the POSIX Threads
programming model. The work was framed around a computational regime rather than
a particular machine: the behaviour of multi-pattern matching when the automaton
greatly exceeds the last-level cache, which is the operating point of
signature-based intrusion detection over realistic rule sets.

The experimental campaign confirmed the principal hypotheses while also
delimiting their scope. Text-level partitioning with a boundary overlap of
$L_{\max}-1$ bytes scales sublinearly, reaching $4.52\times$ at twelve threads
on the multi-gigabyte corpus and a canonical-corpus peak of $4.79\times$ via
dynamic dispatch.
The size of the automaton, and not the content of the corpus, was confirmed to
be the dominant determinant of throughput: as the dictionary grows from a
cache-resident set to one forty-two times larger than the last-level cache,
single-threaded throughput falls by $74.0\%$, and the parallel speedup collapses
accordingly. In the memory-bound regime, the speedup peaks at $2.85\times$ near
six threads (canonical corpus) and then \textit{regresses} as additional threads
are added: the shared DRAM bus saturates before the full thread complement is
engaged, making further concurrency counterproductive.
The corresponding parallel efficiency at twelve threads is $37.7\%$ for the moderate-dictionary regime and $23.3\%$ for the memory-bound one, reflecting DRAM saturation as the dominant overhead rather than thread coordination.
On the heterogeneous processor, frequency-weighted partitioning and
dynamic dispatch were shown to remove the load imbalance of uniform segments,
reducing the run-to-run variation from as much as $53\%$ to under $5\%$; this is a
gain in predictability rather than in peak throughput, since the ceiling remains
the memory subsystem. The proposed flattened output layout, which replaces the
pointer-chasing emission of matches with a contiguous read, proved to be the
cheapest broadly applicable optimization, and its benefit grows with memory
pressure, reaching $+13.0\%$ on a single thread precisely in the memory-bound
regime where it matters most. Parallelizing the construction of the automaton
accelerated the build of the largest dictionary by up to $1.55\times$, a useful
operational gain for rule reloading, though negligible in the total time of a
multi-gigabyte scan.

Equally important are the strategies that were falsified. Partitioning the
dictionary instead of the text yields only a small, locality-driven gain that
saturates early and never approaches the throughput of text parallelism. The
two-dimensional composition of sharding and chunking is consistently slower than
flat-output text chunking, because scanning the text once per shard multiplies the memory
traffic that already constitutes the bottleneck. These negative results are a
substantive part of the contribution: they close plausible regions of the design
space and reinforce the central thesis that, on this class of hardware, the
binding constraint is DRAM bandwidth rather than the matching logic, and that
text parallelism combined with a cache-friendly layout is the response that
composes correctly. End-to-end, the full pipeline achieves up to $4.48\times$ on
the moderate dictionary and $3.04\times$ on the memory-bound one.

The three specific objectives stated in the introduction were each addressed in full.
The first --- to design and implement a parallel Aho-Corasick algorithm in C using
the POSIX Threads API, exploiting intra-input parallelism on shared-memory multi-core
CPUs --- was met through a family of text-partitioning, load-balancing, and
output-layout strategies whose correctness was validated exhaustively across all tested
workloads.
The second --- to quantify performance gains in terms of speedup, throughput, and
parallel efficiency against a sequential baseline --- was met by reporting search
speedups of $4.79\times$ (canonical corpus, dynamic dispatch) and $4.52\times$
(multi-gigabyte corpus, flat-output chunking), alongside parallel efficiencies of
$37.7\%$ and $23.3\%$ for the moderate-dictionary and memory-bound regimes,
respectively.
The third --- to analyse scalability under different workload, pattern-set, and
thread-count configurations --- was met by a systematic sweep spanning seven thread
counts, three corpora, and automata occupying a forty-two-fold range of cache
capacity, from which the memory-bandwidth saturation ceiling and the post-peak
regression were identified and quantified.

\subsecao{Threats to Validity}

The main threat to external validity is the reliance on a single mobile
processor with a constrained thermal envelope, which throttled its frequency
under the sustained, multi-hour measurements and introduced non-trivial variance
in the longer sweeps. The reported figures are deliberately the conservative,
under-load values, and the variability is documented rather than hidden.
Crucially, however, the central conclusion does not depend on the modest cache of
this particular machine: the largest automaton studied ($\sim$507~MiB) exceeds
the last-level cache of nearly every current server processor, and even on the
few parts whose aggregate L3 is large enough to contain it---AMD's stacked
V-Cache models---that capacity is split across core complexes rather than being
uniformly accessible to a single matching thread, so the memory-bandwidth
bottleneck observed here is not an artifact of a small cache but a property that
carries over to server-class hardware. Measuring on a modest chip makes the
phenomenon sharper, while a larger machine with more homogeneous cores and more
memory channels would shift the knee of the curve without removing the wall.
A further limitation is that the microarchitectural explanations, although
consistent with the timing data, were not corroborated with direct hardware
telemetry.

\subsecao{Future Work}

Two directions follow naturally from these limitations. The first is to
instrument the benchmark harness with hardware performance counters, in order to
measure cache misses and pipeline stalls directly and so confirm the
memory-bound explanation with quantitative evidence rather than inference from
end-to-end timing. The second is to repeat the experiments, and in particular the
falsified two-dimensional composition, on server-class processors with much
larger last-level caches and higher memory bandwidth, to locate the cross-over
point at which dictionary sharding might become beneficial. Complementary avenues
include sharding policies driven by the actual state footprint of the
sub-automata rather than by pattern prefixes, and a more rigorous thermal
protocol with controlled cool-down periods to further reduce measurement
variance.
